\section{Términos clave: índice de conceptos de este episodio}

Esta sección reúne los términos de finanzas, IA y mercados de capitales que aparecieron dispersos a lo largo del texto, para que el lector no se quede solo con los nombres de las empresas sin haber asimilado los conceptos. Las explicaciones de la tabla sirven al hilo conductor de este episodio y no sustituyen a textos especializados de finanzas o de aprendizaje automático.

\begin{longtable}{p{0.24\textwidth}p{0.34\textwidth}p{0.33\textwidth}}
\toprule
Término & Explicación en una frase & Papel en este episodio \\
\midrule
Altimeter Capital & Fondo tecnológico de Silicon Valley que invierte en los mercados primario y secundario & Institución de inversión del invitado y origen de su perspectiva. \\
Crossover & Estructura de fondo que invierte tanto en el mercado primario como en el secundario & Explica por qué se puede expresar una tesis sobre la IA a través de Nvidia. \\
CAPEX & Gasto de capital & Inversión de gran escala que requieren por igual los centros de datos de IA, las GPU, el entrenamiento y la reindustrialización. \\
Re-industrialization & Reindustrialización: regreso de la manufactura y la infraestructura a territorio estadounidense & Relacionada con el circuito cerrado de centros de datos, energía, aranceles y capacidad de cómputo para IA. \\
Digitization of Finance & Digitalización de la industria financiera & Contexto común de las stablecoins, los pagos, Robinhood, los mercados de predicción y el Agent Commerce. \\
Stablecoin & Moneda estable: token generalmente anclado a una moneda fiduciaria como el dólar & Posible infraestructura para los pagos de Agents y la digitalización de las finanzas. \\
Negative flywheel & Efecto bola de nieve negativo & Describe la etapa en la que los ingresos de una empresa de modelos no alcanzan a cubrir el costo de entrenar la siguiente generación. \\
Scaling Law & Regularidad empírica entre el tamaño del modelo, los datos, la capacidad de cómputo y el desempeño & Explica por qué las empresas de modelos siguen comprando GPU y entrenando. \\
API & Interfaz para invocar externamente las capacidades de un modelo & Uno de los pilares de ingresos de OpenAI. \\
Agent Commerce & Un Agent busca, compara, compra y paga en nombre del usuario & Podría redefinir los puntos de entrada de la publicidad, el comercio electrónico y los pagos. \\
Alpha & Rendimiento generado por la propia empresa, por encima de la Beta del mercado & La mejora del negocio de Robinhood se usa para ilustrar la Alpha. \\
Beta & Rendimiento que aporta el ciclo de la industria o del mercado & Las criptomonedas, la IA y el mercado tecnológico en su conjunto pueden aportar Beta. \\
Tokenización & Proceso de representar un activo como un token negociable & Una de las vías de Robinhood para entrar a Europa y para la innovación financiera. \\
Prediction Market & Mercado de predicción: usa precios de negociación para expresar la probabilidad de un evento & Se describe como una especie nueva y una muestra de la digitalización de las finanzas. \\
ROIC & Return on Invested Capital, rendimiento sobre el capital invertido & Permite juzgar si la inversión en IA de las grandes empresas tiene retorno. \\
ASIC & Application-Specific Integrated Circuit, chip de propósito específico & Dirección con la que las nubes y las empresas de modelos reducen a largo plazo el riesgo de homogeneización de las GPU. \\
Neo Cloud & Proveedores de servicios en la nube de nueva generación para GPU/IA & Ejemplo de cómo la competencia en la nube se amplía de las tres grandes nubes a más participantes. \\
AI Tax & Propuesta de redistribución para atender los problemas distributivos causados por la sustitución de trabajo por IA & Corresponde al problema social entre el aumento de la productividad y el impacto en el empleo. \\
\bottomrule
\end{longtable}

\subsection{Resumen de la sección}

Los términos de este episodio abarcan finanzas, IA, internet y mercados de capitales. Entender la relación entre ellos permite convertir el programa de «rumores del mundo de la inversión» en una clase sobre estructura industrial: cómo interactúan el gasto de capital, la bolsa de ingresos, el control de los puntos de entrada, la ejecución organizacional y la productividad.
